% !TeX root = ../incremental_SS_Translation.tex

\chapter{Uttaratantra 64: Preservation of Health}% svastharakṣanīya

\section{Literature}

Meulenbeld offered an annotated overview of this chapter, which is called 
\emph{svasthavṛtta} in the vulgate.\fvolcite{IA}[330--331]{meul-hist} There, 
he translates \emph{svasthavṛtta} as ``preservation of health'' but elsewhere 
as ``maintenance of health''\fvolcite{IA}[301]{meul-hist} and 
``hygiene.''\fvolcite{IB}[738]{meul-hist} His only reference to secondary 
literature on this topic is to an annotated German translation by J. Laping 
1984.\footnote{\cite{lapi-1984}, cited by \volcite{IB}[431]{meul-hist}.} His 
footnotes to the chapter also indicate 
many parallels with \emph{Carakasaṃhitā, Sūtrasthāna}, chapter six.

On the term \emph{svastha}, Meulenbeld cites W. Halbfass (1991): 250-252; M. 
Hara (1995) and P.V. Sharma (1985e: 52-53).\fvolcite{IB}[10]{meul-hist}

In primary sources, see the definition of \emph{svasthā} in 
\emph{Suśrutasaṃhitā, Sūtrasthāna} 15.41--44; \emph{svasthahita} in 
\emph{Aṣṭāṅgahṛdaya, Sūtrasthāna} 1.16cd; ...

Comparable chapters in other medical texts include the \emph{svasthacatuṣka} 
in \emph{Carakasaṃhitā, Sūtrasthāna} (see below) and a chapter on 
\emph{svasthavṛtti} in Cakrapāṇidatta's \emph{Cakradatta}.

Meulenbeld also refers to \emph{svasthacatuṣka}, a group of four chapters 
(5--8) in the \emph{Carakasaṃhitā, Sūtrasthāna}, chapters 5–8, also called 
\emph{svasthavṛttacatuṣka} or \emph{svāsthyavṛttikacatuṣka} in various 
colophons of the text and commentaries. He summarises these 
chapters,\fvolcite{IA}[13--14]{meul-hist} and their names are 
\emph{mātrāśitīya}, daily regime (cf. Bhela Sū.8; A.h.Sū.8; A.s.Sū.11. Cf. 
Su.Ci.24); \emph{tasyāśitīya}, seasonal regimen (\emph{ṛtucaryā}) 
(Meulenbeld cites many references in primary and secondary 
literature\fvolcite{IB}[13]{meul-hist}); \emph{navegāndhāraṇīya}, the 
unwholesome effects of suppressing natural urges, etc. (Cf. Bhela Sū.6; 
Su.Ci.24); and \emph{indriyopakramaṇīya}, senses, mind and rules of conduct 
(Cf. Bhela Sū.7; Su.Ci.24. See on Ca.Sū.8: S.Ch. Vidyabhusana (1971): 11--12).

Meulenbeld mentions that a ``Sadānanda wrote a commentary, called \emph{Auṣadhavivṛti}, on the chapters of the
\emph{Carakasaṃhitā} containing rules for the preservation of health 
(svasthavṛtta).\fvolcite{IA}[196]{meul-hist} His footnote on this work refers to 
the NCC (VI, 396–397) and the following bibliographic reference: 
\emph{pañcatantram ... 
carakasūtrasthānasvasthavṛttacatuṣkākhyacaturadhyāyāḥ... 
sadānandaśāstrikṛtauṣadhavivṛtiyutayā saṃvalitaṃ}, Mercantile Press, Lahore 
1926 [IO.San.D.554].\fvolcite{IB}[302]{meul-hist}

In the vulgate, this chapter is included in a section called, ``the chapters that 
are an embellishment of the text" (\emph{tantrabhūṣaṇādhyāya}). However, in 
the Nepalese version, it is the eighth part of tenfold section on treatment for the 
body (\emph{kāyacikitsita}). See K, folio 209r, ll. 4--5 (
\emph{mūtradoṣo mūtrāghāto yonyamānūṣa eva ca} ||
\emph{apasmāronmādakañ caiva rasabhedas tathaiva ca} |
\emph{svastharakṣāvidhāṇañ ca tantrayuktiś ca doṣabhit} |
\emph{ity ebhir daśakaṃ proktaṃ punaḥ kāyacikitsite} ||).


\section{Translation}

\begin{translation}    
  
\item [1] 
Now we shall explain the upkeep of a healthy person. % maintenance?

\item [2] 
One whose humours, digestive fire, and the functioning of bodily constituents and impurities are in equilibrium; whose self, senses, and mind are serene is called healthy.\footnote{This verse is absent here in the vulgate but occurs at the end of \emph{Sūtrasthāna} 15 in both the Nepalese version and vulgate. The \citep[808]{vulgate} supplements its omission with the statement, ``Such a healthy person has been described in the \emph{Sūtrasthāna}'' (\emph{sūtrasthāne samuddiṣṭaḥ svastho bhavati yādṛśaḥ}). }

\item [3] 
O child, the upkeep of that healthy person is the supreme benefit of treatment. The topic of his regime and upkeep, which was mentioned by me at the beginning, was outlined in brief and will here be explained in detail.
% who are the interlocutors?
% something has gone wrong with artham in the Nepalese manuscripts (which has arthāṃ and arthān). The participle (uktaṃ) and verb (vakṣyate) require a singluar. And artham can be neuter. The syntax of the vulgate is wrong because the subject is plural (arthāḥ) and the verb singular (a patch made inadequate by metric causa?)

\item [4] 
In each season, when specific humours are aggravated in people, a wise physician should prescribe the appropriate tastes for that season.

\item [5] 
You should know that when people's digestive fire is weak because of the dampness of their bodies during the rainy months, the bodily winds and so on become aggravated due to their excitement.
% Am I right in thinking that mārutādayaḥ is the five/ten bodily winds and not the doṣas?
% have I over-translated khalu?
    
\end{translation}    
